\subsection{Regular ideals and the class-group correction}
The distinction between the two zeros also changes which multipliers may be used to define ideal classes. Let $R$ be a domain with fraction field $K$, and put $S=R_0$.

\begin{proposition}[The old zero collapses unrestricted ideal classes]\label{prop:class-collapse}
On the ideals $J\ne\{\tau\}$ of $S$, the equivalence relation generated by
\[
 aJ=bJ'\qquad(a,b\in S\setminus\{\tau\})
\]
has one class. If instead one takes nonzero fractional $R$-ideals $I\subset K$, replaces them by $I_0=I\sqcup\{\tau\}\subset K_0$, and permits only multipliers in $K^\times$, the resulting class semigroup is the ordinary fractional ideal class semigroup of $R$. For a Dedekind domain this is its usual ideal class group.
\end{proposition}
\begin{proof}
Every ideal of $S$ distinct from $\{\tau\}$ contains an old ring element, and multiplication by $e=0_R$ then gives
\[
 eJ=\{\tau,e\}.
\]
The allowed choice $a=b=e$ therefore identifies every two such ideals. This proves the first assertion.

For the corrected construction, an $S$-subsemimodule $J$ of $K_0$ containing an element of $K^\times$ is $I\sqcup\{\tau\}$ with $I=J\cap K$. It contains $e$ and is closed under addition, multiplication by $R$, and additive inverses, the last by multiplication by $-1$. Thus $I$ is a nonzero $R$-submodule of $K$. The usual fractional-ideal condition, namely $dI\subset R$ for some $d\in R\setminus\{0\}$, is exactly $dJ\subset S$. Conversely every such $I$ gives an $S$-subsemimodule. Products are the finite sums of pairwise products in both descriptions, so $(IJ)_0=I_0J_0$. Multiplication by $a\in K^\times$ also commutes with this correspondence. Consequently the two class equivalence relations and their products agree. For a Dedekind domain the nonzero fractional ideals are invertible, giving the asserted group; see \citet[Theorem 3.20, p.\ 53]{MilneANT}.
\end{proof}

For a concrete failure of the unrestricted definition, take $R=\Z[\sqrt{-5}]$. The ideal $(2,1+\sqrt{-5})$ has index two: reduction modulo it maps $\sqrt{-5}$ to $1$ in $\F_2$. It is not principal, since a generator would have norm two, but $a^2+5b^2=2$ has no integer solution. Thus its ordinary class is nontrivial, whereas the first equivalence in Proposition~\ref{prop:class-collapse} identifies it with the unit ideal. If regular multipliers are imposed while the old-zero ideal $\{\tau,e\}$ is retained, that ideal contributes an additional absorbing class; restricting the ideals as well is necessary.

\begin{proposition}[Valuation subsemirings]\label{prop:valuations}
Let $V\subset K_0$ be a unital subsemiring containing $\tau$ and $e$. Suppose that for every $a\in K^\times$, at least one of $a,a^{-1}$ belongs to $V$. Then $V=W_0$ for a valuation ring $W$ of $K$. Conversely every valuation ring gives such a subsemiring. For $K=\Q$, these are $\Q_0$ and $(\Z_{(p)})_0$ for rational primes $p$.
\end{proposition}
\begin{proof}
Applying the hypothesis to $-1$ shows that $-1\in V$. Therefore $W=V\cap K$ is a ring and the stated inverse condition is exactly the defining valuation-ring condition. The converse follows from the same definition.

For the rational classification, every subring $W\subset\Q$ containing $\Z$ is a localization of $\Z$ at a set of primes. Indeed, if a reduced fraction $a/b$ lies in $W$, Bezout's identity gives $1/b\in W$, hence the inverses of every prime dividing $b$ lie in $W$. If two distinct primes $p,q$ were not inverted, neither $p/q$ nor $q/p$ would belong to $W$, contradicting the valuation condition. At most one prime is therefore not inverted, giving exactly the rings asserted.
\end{proof}

Finally, the full-support copy of $\A_\Q$ inside $\mathcal A_0$ preserves addition, multiplication, and the multiplicative identity, but sends its additive zero to the old, fully supported zero. It is not a zero-preserving homomorphism of semirings. The zero-adjoined ring $(\A_\Q)_0$ does map into $\mathcal A_0$, with its new zero sent to empty support. This distinction is required whenever the adelic inclusion is used in a categorical assertion.
